% bloch_leads.tex
%
% Purpose:
%   Describe the one-cell scattering matrix, Bloch generalized eigenproblem, and
%   incoming/outgoing trace selection for the periodic half-leads.
%
% Main algorithm:
%   A lead-cell scattering matrix maps incoming Rayleigh amplitudes to outgoing
%   Rayleigh amplitudes.  The Bloch condition turns this map into a generalized
%   eigenproblem whose selected eigenspaces define admissible trace subspaces.
%
% Based on:
%   notes/theory/bloch_mode.md and notes/theory/mode_selection.md.
%
% Main changes:
%   The derivation is written in report form and tied directly to the port
%   trace convention used by the center-cell matrix.
%
% Numerical goal:
%   Explain how D_out^\pm and N_out^\pm are obtained and why they encode the
%   exterior half-lead condition.

\section{Periodic Lead Cells and Bloch Trace Subspaces}

\subsection{One-cell scattering matrix}

For a single lead cell with left wall \(x=X_L\) and right wall \(x=X_R\), write
local Rayleigh amplitudes as \(a^L,b^L,a^R,b^R\), where \(a\) denotes a
right-going local Rayleigh amplitude and \(b\) denotes a left-going local
Rayleigh amplitude.  The finite-dimensional scattering matrix is
\[
  \begin{bmatrix}
  b^L\\
  a^R
  \end{bmatrix}
  =
  S_{\mathrm{cell}}
  \begin{bmatrix}
  a^L\\
  b^R
  \end{bmatrix}
  =
  \begin{bmatrix}
  R_L & T_{R\to L}\\
  T_{L\to R} & R_R
  \end{bmatrix}
  \begin{bmatrix}
  a^L\\
  b^R
  \end{bmatrix}.
\]
In the present code organization, this object is assembled by the Rayleigh
channel builder and scattering-matrix routine in the \texttt{+bloch} package,
together with the boundary integral assembly routines used inside them.

The Bloch condition for a multiplier \(\lambda\) is
\[
  a^R=\lambda a^L,
  \qquad
  b^R=\lambda b^L .
\]
Setting \(a=a^L\) and \(b=b^L\), the scattering relation gives
\[
  b=R_La+\lambda T_{R\to L}b,
  \qquad
  \lambda a=T_{L\to R}a+\lambda R_Rb .
\]
Equivalently,
\[
  \Asc z=\lambda \Bsc z,
  \qquad
  z=\begin{bmatrix}a\\ b\end{bmatrix},
\]
with
\[
  \Asc=
  \begin{bmatrix}
  -R_L&I\\
  T_{L\to R}&0
  \end{bmatrix},
  \qquad
  \Bsc=
  \begin{bmatrix}
  0&T_{R\to L}\\
  I&-R_R
  \end{bmatrix}.
\]
This is the generalized eigenproblem solved by the Bloch-mode module.

\subsection{Trace matrices}

Given a Bloch eigenpair \((\lambda_j,z_j)\), the wall traces are computed in
Rayleigh coordinates.  On the cell left wall,
\[
  D_L(:,j)=a_j^L+b_j^L,
  \qquad
  N_L(:,j)=\ii\GammaR(a_j^L-b_j^L),
\]
where
\[
  \GammaR=\diag(\gamma_m) .
\]
The right wall traces satisfy
\[
  D_R(:,j)=\lambda_jD_L(:,j),
  \qquad
  N_R(:,j)=\lambda_jN_L(:,j).
\]
Here \(N_L,N_R\) denote \(\partial_x u\) trace coefficients on a lead cell wall,
not yet the outward-normal derivatives of the center cell.  The sign conversion
is made when the mode is attached to \(\Gamma_-\) or \(\Gamma_+\).

\subsection{Outgoing trace selection}

For the positive half-lead, the center port \(\Gamma_+\) touches the left wall
of a positive lead cell, and outgoing decay as \(x\to+\infty\) corresponds, away
from unit-circle complications, to
\[
  |\lambda_j|<1 .
\]
Thus
\[
  D_{\mathrm{out}}^+=D_L(:,\mathcal I^+),
  \qquad
  N_{\mathrm{out}}^+=N_L(:,\mathcal I^+),
  \qquad
  \mathcal I^+=\{j:|\lambda_j|<1\}.
\]

For the negative half-lead, the center port \(\Gamma_-\) touches the right wall
of a negative lead cell, and outgoing decay as \(x\to-\infty\) corresponds to
\[
  |\lambda_j|>1 .
\]
Since the center-cell outward normal on \(\Gamma_-\) is \(-e_x\), the Neumann
trace has an additional sign:
\[
  D_{\mathrm{out}}^-=D_R(:,\mathcal I^-),
  \qquad
  N_{\mathrm{out}}^-=-N_R(:,\mathcal I^-),
  \qquad
  \mathcal I^-=\{j:|\lambda_j|>1\}.
\]

When \(k\) is real and propagating Bloch modes are present, some multipliers may
lie on or close to the unit circle.  Then magnitude alone does not decide
incoming versus outgoing.  A typical flux diagnostic is
\[
  \mathcal F_x
  =
  \operatorname{Im}\int_0^d \overline{u(y)}\,\partial_xu(y)\,\dd y
  \approx
  \operatorname{Im}\sum_m \overline{D_m}N_m .
\]
In early numerical experiments it is often more robust to work in a band gap, or
to use \(k+\ii\varepsilon\) with \(\varepsilon>0\), so that the unit-circle
ambiguity is moved away from the selection threshold.

\subsection{What the selected subspace means}

Any arbitrary basis of a \(K\)-dimensional subspace of \(\CC^{2K}\) can close a
finite algebraic system.  It will not, in general, represent the radiation
condition of the periodic half-lead.  The physical content is the subspace
\(\Eout^\pm\) selected by the Bloch problem and the outgoing/admissible rule:
\[
  \Eout^\pm
  =
  \operatorname{range}
  \begin{bmatrix}
  \Dout^\pm\\
  \Nout^\pm
  \end{bmatrix}.
\]
After this subspace is selected, its columns may be orthonormalized or
well-conditioned by QR or SVD without changing the exterior condition.
